\documentclass[11pt]{article}
% Uses the canonical latex-docs semantic style in-repository. The local
% lightweight fallback is for independent compilation of this distribution.
\IfFileExists{technical-implementation.sty}{\usepackage{technical-implementation}}{\usepackage{../support/aaoa-standalone}}
\usepackage{mathtools}
\setDocTitle{Generating Functions and Counting Workbook}
\setDocSubtitle{Guided mathematical exercises, derivations and evidence logs}
\setDocVersion{1.0}
\setDocStatus{Educational reference / unapproved working material}
\setDocOwner{Algorithm Engineering}
\setDocAuthor{Original companion material}
\setDocSummary{Original exercises and analysis space for exact, asymptotic and probabilistic algorithm models.}
\providecommand{\Expect}{\mathbb{E}}
\providecommand{\Prob}{\mathbb{P}}
\providecommand{\Var}{\operatorname{Var}}
\providecommand{\bigoh}{\mathcal{O}}
\begin{document}
\makedoctocpage
\section{Workflow}
State sequence, initial terms, series type, equation, coefficient identity, and finite-$n$ check. For combinatorial objects, first give an unambiguous recursive construction.
\section{Worked calibration}
From $a_n=a_{n-1}+1$ with $a_0=0$, $A(z)=zA(z)+z/(1-z)$ and $A(z)=z/(1-z)^2$. Therefore $a_n=n$.

\section{Exercises and evidence tables}
\subsection{Problem 1: $a_0=2$, $a_n=2a_{n-1}+3$}
Derive OGF and extract coefficients. Test first six values.\par
\textbf{Required evidence:} assumptions, intermediate algebra, boundary check, and a one-sentence model limitation.\par
\vspace{1.2cm}\hrule\vspace{0.4cm}
\subsection{Problem 2: $A(z)=(1-z)^{-4}$}
Derive a finite coefficient formula and the first five terms.\par
\textbf{Required evidence:} assumptions, intermediate algebra, boundary check, and a one-sentence model limitation.\par
\vspace{1.2cm}\hrule\vspace{0.4cm}
\subsection{Problem 3: An ordered binary tree is empty or root-plus-two-ordered-subtrees}
Write a combinatorial class equation, OGF, exact coefficients for $n=0,\ldots,5$, and the limiting scale.\par
\textbf{Required evidence:} assumptions, intermediate algebra, boundary check, and a one-sentence model limitation.\par
\vspace{1.2cm}\hrule\vspace{0.4cm}
\subsection{Problem 4: Four independent Bernoulli$(p)$ trials}
Construct a PGF and use derivatives to obtain mean and variance.\par
\textbf{Required evidence:} assumptions, intermediate algebra, boundary check, and a one-sentence model limitation.\par
\vspace{1.2cm}\hrule\vspace{0.4cm}
\subsection{Problem 5: $F(z,u)$ marks tree size and a parameter}
Explain which coefficient ratio yields the expected parameter and why $u=1$ is necessary.\par
\textbf{Required evidence:} assumptions, intermediate algebra, boundary check, and a one-sentence model limitation.\par
\vspace{1.2cm}\hrule\vspace{0.4cm}
\section{Final verification log}
\begin{longtable}{p{0.12\linewidth}p{0.25\linewidth}p{0.27\linewidth}p{0.24\linewidth}}\toprule Task & Claimed result & Small-instance oracle & Unresolved issue \\ \midrule
1 & & & \\ 2 & & & \\ 3 & & & \\ 4 & & & \\ 5 & & & \\ \bottomrule\end{longtable}
\section{Reflection}
Which derivation depends most strongly on its input law? Identify one assumption whose failure would invalidate the prediction, and define a test that could expose it.

\end{document}
